\section{Design}
\label{sec:design}

\subsection{Agent Metadata}
\AgentMeta contains the request ID, agent ID, role, phase, tool type, graph-node type, prompt-block types, ready time, deadline, run probability, and criticality.
The most important design choice is the fallback key order.
For a request with agent $a$, role $r$, phase $p$, and prompt-block key $b$, \project tries:
\[
(a,r,p,b) \rightarrow (r,p,b) \rightarrow (r,p) \rightarrow (r) \rightarrow global.
\]
Specific keys capture stable agent behavior when enough samples exist.
Coarser keys avoid cold-start failure.

\subsection{Route Signatures}
\RouteSig maintains expert counts per key and layer.
For each key-layer pair it derives a probability distribution, top-$M$ experts, entropy, confidence, miss costs, sample count, and update timestamp.
Confidence combines sample count, distribution concentration, and recent active-set stability.
High-confidence signatures can drive aggressive prefetch; low-confidence signatures can be used only for soft scheduling or ignored.

\begin{algorithm}[t]
\caption{\RouteSig lookup}
\label{alg:routesig}
\begin{algorithmic}[1]
\Procedure{Lookup}{$meta, layer, min\_confidence$}
  \For{$key$ in \Call{FallbackKeys}{$meta$}}
    \State $sig \gets$ \Call{Signature}{$key, layer$}
    \If{$sig \neq \emptyset$ and $sig.confidence \geq min\_confidence$}
      \State \Return $sig$
    \EndIf
  \EndFor
  \State \Return $\emptyset$
\EndProcedure
\end{algorithmic}
\end{algorithm}

\subsection{Serving Policies}
\project evaluates three policies.
\textbf{Prefetch} loads the predicted top-$M$ experts before router execution and counts hits, misses, wasted prefetches, and a stall-reduction proxy.
\textbf{Scheduling} groups requests whose predicted \EWSs overlap, while preserving dependency order.
\textbf{Placement} uses predicted demand to evaluate an EPLB-style expert movement decision, accepting moves that reduce tail load and rejecting unstable moves.
All three policies are advisory: exact routing remains authoritative.

\subsection{Capability Gate}
The vLLM sidecar inspects the model config before requesting routed-expert capture.
Fields such as \texttt{num\_local\_experts}, \texttt{num\_experts\_per\_tok}, or architecture names containing ``MoE'' enable routed capture.
Dense configs return a structured fallback reason.
This design keeps one serving path for dense and MoE models while making unsupported optimization explicit.
